
This annex describes how \openshmem programs use the GPU extensions defined in
this specification.

A typical \openshmem GPU program proceeds as follows:

\begin{enumerate}
\item The desired \ac{GPU} is selected and made current using the vendor-specific
  GPU runtime (e.g., \texttt{cudaSetDevice}, \texttt{hipSetDevice}).
\item The \openshmem library is initialized with GPU support by calling
  \FUNC{shmem\_init\_thread} with \CONST{SHMEM\_DEVICE\_HOST\_INIT},
  \CONST{SHMEM\_DEVICE\_KERNEL\_INIT}, or both.
\item The \VAR{provided} argument is examined to confirm that the requested
  capabilities are available.
\item Symmetric memory is allocated from the \ac{GPU} symmetric heap using
  \FUNC{shmemg\_malloc}, \FUNC{shmemg\_calloc}, or \FUNC{shmemg\_align}.
\item Communication is performed:
  \begin{itemize}
  \item \textbf{GPU-aware:} standard \openshmem routines are invoked from the
    host with GPU-attached memory as operands.
  \item \textbf{GPU-centric:} device kernels that invoke \texttt{shmemg\_}
    routines directly are launched.
  \end{itemize}
\item \ac{GPU} symmetric heap memory is freed using \FUNC{shmemg\_free}.
\item The \openshmem library is finalized using \FUNC{shmem\_finalize}.
\end{enumerate}

\subsection*{Example: GPU-Aware Put}

\lstinputlisting[style=SourceListingsDefault]{example_code/gpu_aware_put_example.c}

\subsection*{Example: GPU-Centric Put}

\lstinputlisting[style=SourceListingsDefault]{example_code/gpu_centric_put_example.c}

\clearpage

